% ===== 原文 pp.113--118 =====
% [原文 p.113]
\begin{figure}[htbp]
\centering
\includegraphics[width=0.86\textwidth]{fig25.png}
\figtitle{25}{用另一个拓扑边界条件 $\mathfrak{B}^{\mathrm{sym}}_{\cS'}$ 施行夹心构造，得到理论 $\cT'$，它由规范化（部分）对称性 $\cS$ 得到，可能带有离散挠率。$\mathfrak{B}^{\mathrm{sym}}_{\cS'}$ 的拓扑算符生成 $\cT'$ 的对称性 $\cS'$，后者是规范化下对偶对称性与剩余对称性的组合。}
\end{figure}

现在用边界条件 $\mathfrak{B}^{\mathrm{sym}}_{\cS'}$ 施行夹心构造，得到一个具有对称性 $\cS'$ 的 $d$ 维理论 $\cT'$。这两个 $d$ 维理论由一个离散规范化相联系。也就是说，我们可以通过规范化 $\cT$ 的对称性 $\cS$ 的一部分（可能带有离散挠率）来得到 $\cT'$。见图 25。$\cS'$ 是由规范化产生的对偶对称性与未被规范化的剩余对称性组合而成的对称性。

\begin{examplebox}{例 4.18：规范化一个高阶形式对称性}
考虑一个理论 $\cT$，它带有一个无反常的离散 $p$-形式对称性 $G^{(p)}$（当 $p=0$ 时我们只考虑阿贝尔的 $G^{(0)}$）。如上文所讨论，SymTFT 的作用量为
\begin{equation}
  2\pi\int a_{p+1}\cup_\eta\delta b_{d-p-1}\,. \tag{4.261}
\end{equation}

设我们选取如下的拓扑边界条件 $\mathfrak{B}^{\mathrm{sym}}_{\cS'}$：
\begin{itemize}[leftmargin=1.4em,itemsep=1pt,topsep=1pt]
\item 对规范场 $a_{p+1}$ 取 \textbf{Neumann} 边界条件（b.c.）；
\item 对规范场 $b_{d-p-1}$ 取 \textbf{Dirichlet} 边界条件。
\end{itemize}

此时夹心构造给出的理论是
\begin{equation}
  \cT'=\cT/G^{(p)}\,. \tag{4.262}
\end{equation}

的确，$\mathfrak{B}^{\mathrm{sym}}_{\cS'}$ 上的拓扑算符是 $U_{\hat g}$，它们生成
\begin{equation}
  G^{(d-p-2)}=\widehat{G}^{(p)}\,, \tag{4.263}
\end{equation}
% 原文方框跨 p.113--p.114，此处按 STYLE_GUIDE §7 合并为一个方框，下面承接原书 p.114 页首的收尾句。
% [原文 p.114]
即 $\cT/G^{(p)}$ 的对称性。拓扑算符 $U_g$ 现在终止在 $\mathfrak{B}^{\mathrm{sym}}_{\cS'}$ 上，从而给出 $\cT/G^{(p)}$ 的带电算符。
\end{examplebox}

% [Tutor Note: p.113--p.114 的例 4.18 在原文中是一个跨页方框，此处合并为一个 `examplebox`，编号 (4.261)--(4.263) 只出现一次、未被拆成两个。核对 (4.263)：规范化某 $p$-形式对称性 $G^{(p)}$ 后得到的对偶对称性是 $(d-p-2)$-形式，在 $d$ 维时空中与 $p$-形式算符所在余维互补，且对偶群为 $\widehat{G}^{(p)}$，与本文 4.3.3 节的结论一致。]

\para{规范化的障碍。}规范化的障碍，即 $\cS$ 的 't Hooft 反常，现在表现为 $\cZ(\cS)$ 的拓扑边界条件存在的障碍。再次考虑高阶形式对称性，此时 SymTFT 的作用量为 (4.258)。$\cL^{\mathrm{eff}}_{d+1}$ 的存在把 $b_{d-p_i-1}$ 场的运动方程修改为
\begin{equation}
  \delta b_{d-p_i-1}=f_i[a_{p_1+1},a_{p_2+1},\cdots,a_{p_n+1}]\,, \tag{4.264}
\end{equation}
其中 $f_i$ 是某个 $(d-p_i)$-上链。现在，对某个固定的 $i$，我们只有在同时对足够多的 $a_{p_j+1}$ 场取 Dirichlet 边界条件、使得 $f_i$ 完全成为背景场时，才能对 $b_{d-p_i-1}$ 取 Dirichlet 边界条件。这个条件未必可解，因为它可能要求对同一个 $i$ 同时对 $b_{d-p_i-1}$ 与 $a_{p_i+1}$ 取 Dirichlet 边界条件，而这是矛盾的。让我们用一个例子明确地看清这一点：

\begin{examplebox}{例 4.19：纯 $\SO(6)$ Yang--Mills 理论}
回忆例 4.8：纯 $\SO(6)$ 规范理论在 $G^{(1)}=\Z_2^e$ 电 1-形式对称性与 $G^{(d-3)}=\Z_2^m$ $(d-3)$-形式对称性之间有混合 't Hooft 反常，由
\begin{equation}
  \cI_{d+1}[A_2^e,A_{d-2}^m]=\exp\left(\pi\ii\int A_{d-2}^m\cup\operatorname{Bock}(A_2^e)\right). \tag{4.265}
\end{equation}
给出。SymTFT 的作用量为
\begin{equation}
  \pi\int\left(a_2^e\cup\delta b_{d-2}^e+a_{d-2}^m\cup\delta b_2^m+a_{d-2}^m\cup\operatorname{Bock}(a_2^e)\right). \tag{4.266}
\end{equation}
由 $a_{d-2}^m$ 导出的运动方程是
\begin{equation}
  \delta b_2^m=\operatorname{Bock}(a_2^e)\,, \tag{4.267}
\end{equation}
这意味着对 $b_2^m$ 取 Dirichlet 边界条件要求对 $a_2^e$ 取 Dirichlet 边界条件。换言之，我们不能同时对 $b_2^m$ 与 $b_{d-2}^e$ 取 Dirichlet 边界条件，因为后一 Dirichlet 边界条件意味着对 $a_2^e$ 取 Neumann 边界条件。换个说法，我们不能同时规范化电与磁的高阶形式对称性，这也正是 't Hooft 反常 (4.265) 所施加的障碍。
\end{examplebox}

% [Tutor Note: (4.265) 的系数是 $\pi\ii$（不是 $2\pi\ii$）：因为 $\Z_2$ 上链取值为 $0,\frac12$，$\pi\ii\int A^m\cup\operatorname{Bock}(A^e)$ 给出 $\pm1$ 的相因子，正是 $\Z_2$ 反常的标准归一化。校验 (4.266)：$a^m_{d-2}$ 的运动方程来自 $\delta(a^m_{d-2}\cup\delta b^m_2+a^m_{d-2}\cup\operatorname{Bock}(a^e_2))$，对 $b^m_2$ 变分得 $\delta b^m_2=\operatorname{Bock}(a^e_2)$，与 (4.267) 自洽。]

\section{高阶群对称性}

当不同次数的高阶形式对称性彼此混合时，就出现高阶群对称性。这意味着：对 $p$-形式对称性的背景场 $B_{p+1}$ 作一次规范变换，
% [原文 p.115]
也会移动 $q$-形式对称性的背景场 $B_{q+1}$。若 $r$ 是参与该高阶群的 $r$-形式对称性的最大整数，则该高阶群是一个 $(r+1)$-群。我们将主要关注最简单类型的高阶群，即只涉及 0-形式与 1-形式对称性的 2-群对称性。

% [Tutor Note: §5 的「(r+1)-群」命名规则：参与混合的对称性最高形式次数为 $r$，则群称为 $(r+1)$-群；这与「$n$-群」的一般定义（Postnikov 塔在 $n+1$ 层截断）一致。核对：只含 0-形式与 1-形式时为 $r=1$，故称 2-群，与 5.2.2 的 2-群定义相符。]

关于这一课题已有大量工作。讨论 QFT 中高阶群对称性的奠基性论文见 [181--184]。通过弦论以及高维 QFT 的紧致化来讨论强耦合 QFT 中高阶群对称性的工作始于 [185--187]。关于高阶群对称性，其他一些工作（排名不分先后）的样本见 [188--206]。

\subsection{连续高阶群对称性}

\para{例：规范化混合 't Hooft 反常。}让我们从一个高阶群对称性的例子开始，其中各个组成的 $p$-形式对称群都是连续的 [183]。这类连续高阶群的已规范化版本在反常抵消的 Green--Schwarz 机制中著名地出现。

让我们从一个 $d$ 维理论出发，它具有一个 $p$-形式对称性 $\U(1)^{(p)}$ 与一个 $(d-2p-4)$-形式对称性 $\U(1)^{(d-2p-4)}$，二者带有混合 't Hooft 反常
\begin{equation}
  W[B_{p+1}+\dd\Lambda_p,B_{d-2p-3}+\dd\Lambda_{d-2p-4}]
  =W[B_{p+1},B_{d-2p-3}]+2\pi\kappa\int\Lambda_p\wedge H_{p+2}\wedge H_{d-2p-2}\,. \tag{5.1}
\end{equation}
这里 $\kappa\in\Z$ 刻画该反常，$W$ 是有效作用量，$H_{p+2}$ 与 $H_{d-2p-2}$ 分别是背景场 $B_{p+1}$ 与 $B_{d-2p-3}$ 的场强。

由于不存在依赖于 $\Lambda_{d-2p-4}$ 的反常项，我们可以规范化 $\U(1)^{(d-2p-4)}$ 对称性。天真地看，规范化会把该混合 't Hooft 反常转变成一个 ABJ 反常，这会表明规范化后的理论不再带有 $\U(1)^{(p)}$ 对称性。然而，这个 ABJ 反常可以按类似于 Green--Schwarz 机制的方式抵消。注意，规范化过程引入一个新的 $(2p+1)$-形式对称性 $\U(1)^{(2p+1)}$，它由 Noether 流
\begin{equation}
  j^{d-2p-2}=H_{d-2p-2}\,. \tag{5.2}
\end{equation}
生成。如我们所知，该对称性的背景场 $B_{2p+2}$ 以下式耦合：
\begin{equation}
  2\pi\int B_{2p+2}\wedge H_{d-2p-2}\,. \tag{5.3}
\end{equation}
天真地看，在 $\U(1)^{(2p+1)}$ 规范变换下，这个背景场按 $B_{2p+2}\to B_{2p+2}+\dd\Lambda_{2p+1}$ 变换；但如果我们修改 $B_{2p+2}$ 的规范变换，使得 $B_{2p+2}$ 也在 $\Lambda_p$ 下变换，我们就能通过在作用量中增添 (5.3) 这个新项来抵消 ABJ 反常 (5.1)。

% [Tutor Note: 式 (5.1) 右端第一项是 $W[B_{p+1},B_{d-2p-3}]$ 的规范变换「原样」，反常是第二项 $2\pi\kappa\int\Lambda_p\wedge H_{p+2}\wedge H_{d-2p-2}$；注意被积的 $\Lambda_p$ 是 $p$-形式规范参数，次数核对：$p+(p+2)+(d-2p-2)=d$，故该项是 $d$-形式，恰是 $d$ 维有效作用量 $W$ 在一次规范变换下产生的相位（不是 $(d+1)$-形式）。把同一反常写成 $(d+1)$ 维反常理论/SymTFT 作用量的形式见 §4.2.3 与 (4.197)。]

% [原文 p.116]
更明确地说，令 $B_{2p+2}$ 按如下方式变换：
\begin{equation}
  B_{2p+2}\to B_{2p+2}+\dd\Lambda_{2p+1}-\kappa\Lambda_p\wedge H_{p+2}\,. \tag{5.4}
\end{equation}
那么反常 (5.1) 被抵消，且有效作用量按如下方式变换：
\begin{equation}
  W[B_{p+1}+\dd\Lambda_p,B_{2p+2}+\dd\Lambda_{2p+1}-\kappa\Lambda_p\wedge H_{p+2}]
  =W[B_{p+1},B_{2p+2}]\,. \tag{5.5}
\end{equation}
$B_{2p+2}$ 在 $B_{p+1}$ 的规范变换下发生变换这一事实，意味着 $\U(1)^{(p)}$ 与 $\U(1)^{(2p+1)}$ 对称性组合成一个 $(2p+2)$-群对称性。最简单的情形是 $p=0$，此时我们得到一个 \textbf{2-群}对称性。

注意，通常的场强 $H_{2p+3}=\dd B_{2p+2}$ 不再是规范不变的，而需要修改为
\begin{equation}
  H_{2p+3}=\dd B_{2p+2}+\kappa B_{p+1}\wedge H_{p+2}\,, \tag{5.6}
\end{equation}
它现在不再闭合。事实上，刻画高阶群、含 $\kappa$ 的项，恰好被 $H_{2p+3}$ 闭合性的破坏所捕捉：
\begin{equation}
  \dd H_{2p+3}=\kappa H_{p+2}\wedge H_{p+2}\,. \tag{5.7}
\end{equation}

\para{一般结构。}一般地，若对某个连续 $p$-形式对称性的场强 $H_{p+2}$ 有关系式
\begin{equation}
  H_{p+2}=\dd B_{p+1}+\Theta(B_{p_1+1},B_{p_2+1},\cdots,B_{p_n+1})\,, \tag{5.8}
\end{equation}
其中 $\Theta$ 是其他连续 $p_i$-形式对称性的背景场 $B_{p_i+1}$ 的函数，且使得 $H_{p+2}$ 在所有规范变换下不变，则我们说存在一个连续高阶群对称性。

% [Tutor Note: 由 (5.6) 验证 (5.7)：$\dd H_{2p+3}=\dd^2B_{2p+2}+\kappa\,\dd B_{p+1}\wedge H_{p+2}+\kappa B_{p+1}\wedge \dd H_{p+2}$；因 $\dd H_{p+2}=0$、$\dd^2=0$，只剩 $\kappa H_{p+2}\wedge H_{p+2}$。对 $p=0$ 情形这就是 Bockstein 2-群的连续对应物。]

\subsection{离散高阶群对称性}

\subsubsection{动机：群扩张}

离散高阶群对称性涉及用其他 $p_i$-形式对称性对 $p$-形式对称性的扩张。因此我们先回顾最简单的扩张情形，即阿贝尔群的扩张。

\para{群扩张的扩张类。}设 $K^{(0)}$ 与 $H^{(0)}$ 是两个阿贝尔群，它们处于短正合列
\begin{equation}
  0\to H^{(0)}\to G^{(0)}\to K^{(0)}\to 0\,. \tag{5.9}
\end{equation}

% [原文 p.117]
若该短正合列不分裂，则我们不能把 $G^{(0)}$ 表成 $H^{(0)}$ 与 $K^{(0)}$ 的直积
\begin{equation}
  G^{(0)}\ncong H^{(0)}\times K^{(0)}\,. \tag{5.10}
\end{equation}
物理上，我们应把这些群看作理论的 0-形式对称性。

理解扩张 $G^{(0)}$ 的一种方式是把其元素表示为对
\begin{equation}
  g=(h,k)\,,\qquad h\in H^{(0)},\ k\in K^{(0)}\,. \tag{5.11}
\end{equation}
于是群运算取形式
\begin{equation}
  g_1+g_2=(h_1+h_2+e(k_1,k_2),k_1+k_2)\,, \tag{5.12}
\end{equation}
它是用 $K^{(0)}$ 上取值于 $H^{(0)}$ 的 2-上链 $e$ 来描述的。要求群运算满足结合律，意味着 $e$ 是一个 2-上闭链
\begin{equation}
  \delta e=0\,. \tag{5.13}
\end{equation}
最后还有一个平移的自由度
\begin{equation}
  (h,k)\to(h+f(k),k)\,, \tag{5.14}
\end{equation}
它不改变扩张，其中 $f$ 是 $K^{(0)}$ 上取值于 $H^{(0)}$ 的 1-上链。这一平移把 $e$ 修改为
\begin{equation}
  e\to e+\delta f\,. \tag{5.15}
\end{equation}
于是，短正合列 (5.9)（因而扩张群 $G^{(0)}$ 的形式）由一个上同调类
\begin{equation}
  [e]\in H^2(K^{(0)},H^{(0)}) \tag{5.16}
\end{equation}
描述。

% [Tutor Note: (5.12)--(5.16) 是 Eilenberg--MacLane 群扩张理论的标准结论：$H^2(K^{(0)},H^{(0)})$ 分类 $K^{(0)}$ 被 $H^{(0)}$ 的中心扩张。(5.13) 的 $\delta$ 是群上链微分，(5.15) 表示 $e$ 的平移正是上同调等价。校验：令 $e=0$ 即得直积群，对应 $H^2$ 的零类。]

\para{用拓扑算符解释。}上闭链 $e$ 的物理含义如下：当两个与 $k_1,k_2\in K^{(0)}$ 相应的余维-1 拓扑算符合到一起、形成与 $k_1k_2\in K^{(0)}$ 相应的拓扑算符时，所得的余维-2 联结充当与
\begin{equation}
  e(k_1,k_2)\in H^{(0)} \tag{5.17}
\end{equation}
相应的余维-1 拓扑算符的源。见图 26。

% [原文 p.118]
\begin{figure}[htbp]
\centering
\includegraphics[width=0.6\textwidth]{fig26.png}
\figtitle{26}{用 $K^{(0)}$（蓝色）与 $H^{(0)}$（红色）的拓扑余维-1 算符表示群扩张 (5.9)。}
\end{figure}

\para{对背景场的含义。}$K^{(0)}$ 拓扑算符的联结为 $H^{(0)}$ 拓扑算符提供源，这一事实意味着：当开启 $K^{(0)}$ 背景场时，$H^{(0)}$ 的背景场不再闭合。容易看出，这两个背景场的关系为
\begin{equation}
  \delta B_1^H=\left(B_1^K\right)^*e\,, \tag{5.18}
\end{equation}
其中 $B_1^H$ 与 $B_1^K$ 分别是 $H^{(0)}$ 与 $K^{(0)}$ 的背景场，而 $\left(B_1^K\right)^*e$ 是 $e$ 通过 $B_1^K$ 的拉回，这类拉回曾在 (4.76) 之后讨论过。

% [Tutor Note: (5.18) 中 $\left(B_1^K\right)^*e$ 的记号：$B_1^K$ 作为映射 $M\to BK^{(0)}$，$e$ 是 $K^{(0)}$ 上的 2-上链，拉回后成为时空上的 2-上链；$\delta B_1^H$ 是 2-上链，$(B_1^K)^*e$ 也是 2-上链，两边次数相同；这里用群上链微分 $\delta$ 而非 $\dd$，是因为背景场是离散群取值的上链。核对：图 26 中红线的端点（源）由蓝线交点的余维-2 联结给出。]

\subsubsection{一般结构}

\para{高阶群。}推广 (5.18)，我们说存在一个离散高阶群对称性，如果某个 $G^{(p)}$ $p$-形式对称性的背景场不闭合
\begin{equation}
  \delta B_{p+1}+\Theta(B_{p_1+1},B_{p_2+1},\cdots B_{p_n+1})=0\,, \tag{5.19}
\end{equation}
其中 $\Theta$ 是时空上取值于 $G^{(p)}$ 的 $(p+2)$-上闭链，它是其他背景场 $B_{p_i+1}$ 的函数。类似地，高阶群中其他高阶形式对称性的背景场也可能不闭合。

(5.19) 有一个漂亮的拓扑算符解释：$\Theta(B_{p_1+1},B_{p_2+1},\cdots B_{p_n+1})$ 的 Poincaré 对偶是一个 $(d-p-2)$-圈，它充当余维-$(p+1)$ 的 $G^{(p)}$ 拓扑算符的源。

\para{2-群。}让我们专门讨论由 $G^{(0)}$ 0-形式对称性与 $G^{(1)}$ 1-形式对称性构成的 2-群对称性。2-群对称性由一个 \textbf{Postnikov 类}
\begin{equation}
  [\Theta]\in H^3\left(G^{(0)},G^{(1)}\right), \tag{5.20}
\end{equation}
指定。

% [Tutor Note: (5.19) 的次数校验：$\Theta$ 是 $(p+2)$-上闭链而 $\delta B_{p+1}$ 是 $(p+2)$-上链，故等式两边次数匹配。在 (5.20) 中取 $p=1$ 得 $\Theta$ 为 3-上闭链，与 $H^3(G^{(0)},G^{(1)})$ 的上标 3 相符；这是「2-群 ↔ $H^3$」的来历，也是把 Postnikov 类称为「3-类」的原因。]
